\documentclass[11pt]{article}

% ---------------------------------------------------------------------------
%  The high-Reynolds tau-saturation residual floor of the stabilized porous
%  Navier-Stokes method, contrasted with pure Galerkin.
%
%  Companion note to theory/paper/article.tex
%  ("A stabilized finite element method for incompressible, inertial flows in
%   inhomogeneous porous media", Casas, Gonzalez-Usua, Codina, de-Pouplana).
%  ALL notation, operators, stabilization parameters and asymptotic regimes
%  are taken verbatim from that paper; cross-references of the form
%  (paper, eq:Label) point to its equation labels.
%
%  Compiles with a standard TeX install:  pdflatex tau_saturation_note.tex
% ---------------------------------------------------------------------------

\usepackage[margin=1in]{geometry}
\usepackage{amsmath,amssymb,amsthm}
\usepackage{bm}
\usepackage{booktabs}
\usepackage{hyperref}

% --- notation, matched to theory/paper/article.tex ------------------------
\newcommand{\bu}{\boldsymbol{u}}
\newcommand{\buh}{\boldsymbol{u}_h}
\newcommand{\bw}{\boldsymbol{w}}
\newcommand{\ba}{\boldsymbol{a}}
\newcommand{\bv}{\boldsymbol{v}}
\newcommand{\bvh}{\boldsymbol{v}_h}
\newcommand{\bef}{\boldsymbol{f}}
\newcommand{\bn}{\boldsymbol{n}}
\newcommand{\bx}{\boldsymbol{x}}
\newcommand{\bsig}{\boldsymbol{\sigma}}
\newcommand{\ViscProj}{\overset{\scriptscriptstyle{\mathrm{DS}}}{\Pi}} % paper's Xi
\newcommand{\Lop}{\mathcal{L}}
\newcommand{\Ladj}{\mathcal{L}^{*}}
\newcommand{\Rop}{\mathcal{R}}
\newcommand{\tauK}{\boldsymbol{\tau}_K}
\newcommand{\tauNS}{\tau_{1,\mathrm{NS}}}
\newcommand{\sigtil}{\widetilde{\sigma}_{\alpha}}
\newcommand{\Reyn}{\mathrm{Re}}
\newcommand{\Damk}{\mathrm{Da}}
\newcommand{\BS}{B_{\mathrm S}}
\newcommand{\LS}{L_{\mathrm S}}
\newcommand{\rG}{r_{\mathrm{G}}}
\newcommand{\rS}{r_{\mathrm{S}}}
\newcommand{\rM}{r_{\mathrm{M}}}
\newcommand{\dd}{\,\mathrm{d}}
\newcommand{\grad}{\nabla}
\newcommand{\divg}{\nabla\!\cdot\!}
\newcommand{\bigOh}{\mathcal{O}}
\newcommand{\innerK}[2]{\langle #1, #2\rangle_K}
\newcommand{\normh}[1]{\big\|#1\big\|_h}
\newcommand{\triplenorm}[1]{{\lvert\kern-0.25ex\lvert\kern-0.25ex\lvert #1
  \rvert\kern-0.25ex\rvert\kern-0.25ex\rvert}}

\newtheorem{proposition}{Proposition}
\newtheorem{remark}{Remark}

\title{\bf The high-Reynolds $\tau_1$-saturation residual floor\\
of the stabilized porous Navier--Stokes method, contrasted with Galerkin}
\author{Companion note to \texttt{article.tex} --- porous Navier--Stokes VMS solver}
\date{\today}

\begin{document}
\maketitle

\begin{abstract}
This note explains, in the notation and within the analytical framework of the
paper \texttt{article.tex}, why the equal-order stabilized (ASGS/OSGS) method
develops an irreducible \emph{algebraic residual floor} at high Reynolds number
when judged by a tight, scale-free, \emph{absolute} stopping gate, whereas the
pure Galerkin discretization does not. The mechanism is the saturation of the
stabilization parameter $\tau_1$ defined in (paper, \texttt{eq:Tau1Final}) as the
viscosity $\nu\to0$ -- the same saturation that underlies the paper's
dominant-convection asymptotics (paper, \S\,\texttt{sec:Robustness},
$\tau_1\sim h/(\alpha\lVert\ba\rVert_{\infty,K})$) and the paper's own warning,
below the Galerkin coercivity identity (paper, Eq.~after \texttt{eq:Galerkin
coercivity}), that ``for very small viscosities or fluid volume fractions [the
viscous term] will provide almost no control over the gradient of the velocity.''
We make two scopes explicit, both consistent with the paper: the floor is
\emph{generic} to high-$\Reyn$ Navier--Stokes (it survives $\bsig\equiv0$), and
the additional difficulty of the low-porosity cells is the \emph{separate},
ubiquitous $1/\alpha_0$ degradation that the paper's robustness analysis already
predicts. Crucially, the paper's convergence theorem (paper,
\texttt{eq:ConvergenceResult}) guarantees the solution error stays optimal, so
the floor is a stopping-criterion artifact, not an under-resolved solution.
\end{abstract}

\section{The problem, operators and parameters (paper notation)}\label{sec:setup}

We reproduce only what is needed; every object is that of \texttt{article.tex}.
On $\Omega\subset\mathbb{R}^d$ with fluid volume fraction $\alpha:\Omega\to(0,1]$,
kinematic viscosity $\nu$, and symmetric positive-semidefinite resistance tensor
$\bsig(\alpha,\bu)$, the strong porous Navier--Stokes problem is (paper,
\texttt{eq:StrongMomentumEquation}--\texttt{eq:StrongMassEquation})
\begin{equation}
  \alpha\,\bu\!\cdot\!\grad\bu \;-\; 2\divg(\alpha\nu\,\ViscProj\grad\bu)
  \;+\;\alpha\grad p \;+\;\bsig(\alpha,\bu)\bu \;=\;\bef,
  \qquad
  \varepsilon p + \divg(\alpha\bu)=0,
  \label{eq:strong}
\end{equation}
with $\ViscProj$ the deviatoric--symmetric projection of the velocity gradient and
$\varepsilon\ge0$ the iterative-penalty compressibility. The Darcy--Brinkman--%
Forchheimer closure is $\bsig=\sigma(\alpha,\bu)\,\mathbb I$ with (paper,
\texttt{eq:DBFResistanceTerm})
\begin{equation}
  \sigma(\alpha,\bu)=a(\alpha)+b(\alpha)\,|\bu| .
  \label{eq:dbf}
\end{equation}
Writing $U=[\bu;p]$, $F=[\bef;0]$, the problem is cast (paper,
\texttt{eq:CDR\_equation}) as $\Lop_{\bu}U=F$ with the convection--diffusion--%
reaction operator $\Lop_{\bw}U=-\partial_i(\mathbf K_{ij}\partial_jU)
+\mathbf A_{\mathrm c,i}(\bw)\partial_iU+\mathbf A_{\mathrm f,i}\partial_iU
+\mathbf S(\bw)U$, and \emph{residual operator}
$\Rop_{\bw}U_h := F-\Lop_{\bw}U_h$.

\paragraph{Galerkin and stabilized forms.}
The Galerkin bilinear form is (paper, \texttt{eq:823})
\begin{equation}
  B(\ba,U_h,V_h)=\langle\bvh,\alpha\,\ba\!\cdot\!\grad\buh\rangle
  +2\langle\grad\bvh,\alpha\nu\ViscProj\grad\buh\rangle
  +\langle\bvh,\alpha\grad p_h\rangle
  +\langle\bvh,\sigma\buh\rangle
  +\langle q_h,\varepsilon p_h\rangle
  +\langle q_h,\divg(\alpha\buh)\rangle .
  \label{eq:B}
\end{equation}
The VMS stabilized problem (paper, \texttt{eq:OSGSProblem}) reads
$\BS(\bw,U_h,V_h)=\LS(\bw,V_h,\boldsymbol\pi_h)$ with (paper, Eqs.~under
\texttt{eq:OSGSProblem})
\begin{equation}
  \BS(\bw,U,V)=B(\bw,U,V)-\sum_K\innerK{\Ladj_{\bw}V}{\tauK(\bw)\,\Lop_{\bw}U},
  \quad
  \LS(\bw,V,\boldsymbol\pi)=L(V)-\sum_K\innerK{\Ladj_{\bw}V}{\tauK(\bw)\,(F-\boldsymbol\pi)},
  \label{eq:BS}
\end{equation}
where $\Ladj_{\bw}$ is the adjoint operator (paper, \texttt{eq:AdjointDifferential%
Operator}); ASGS is $\boldsymbol\pi_h=\boldsymbol0$, OSGS is
$\boldsymbol\pi_h=\Pi_{\boldsymbol\tau h}\Rop U_h$. (As in the paper, the term
$\tfrac1\alpha\divg(\alpha\ba)\bvh$ is dropped from $\Ladj V_h$ to preserve the
symmetry of the stability estimate.)

\paragraph{Stabilization parameters.}
With the porosity well resolved by the mesh and $\bsig=\sigma\mathbb I$ uniform,
the parameters are (paper, \texttt{eq:Tau1Final}, \texttt{eq:Tau2Final},
\texttt{eq:TauNavierStokes})
\begin{equation}
  \boxed{\;\tau_1=\frac{1}{\alpha_K\,\tauNS^{-1}+\sigma}\;},
  \qquad
  \tau_2=\frac{h^2}{c_1\alpha_K\tauNS},
  \qquad
  \tauNS=\Big(c_1\frac{\nu}{h^2}+c_2\frac{|\bw|}{h}\Big)^{-1},
  \label{eq:tau}
\end{equation}
with $\alpha_K$ a representative (taken as the maximal) porosity on $K$, $|\bw|$
the elemental velocity modulus, and, for equal-order order-$k$ interpolation,
$c_1=4k^4$, $c_2=2k^2$ (paper, Remark after \texttt{eq:conditions\_on\_num\_param}).

\section{The Galerkin coercivity already loses control as $\nu\to0$}\label{sec:galcoerc}

Testing \eqref{eq:B} with $U_h$ itself, under $\divg(\alpha\ba)=0$ and homogeneous
Dirichlet data, gives the Galerkin coercivity identity (paper, \texttt{eq:Galerkin
coercivity})
\begin{equation}
  B(\ba,U_h,U_h)=2\nu\big\|\alpha^{1/2}\ViscProj\grad\buh\big\|^2
  +\big\|\sigma^{1/2}\buh\big\|^2+\varepsilon\|p_h\|^2 .
  \label{eq:galcoerc}
\end{equation}
The paper observes immediately below this identity that \emph{``for very small
viscosities or fluid volume fractions, the first term above will provide almost no
control over the gradient of the velocity, leading to oscillations on the
solution; this is what happens in the standard case, but here the problem is
aggravated for small porosities.''} This is the root cause, stated at the
continuous-analysis level: as $\nu\to0$ the only velocity-gradient control in
\eqref{eq:galcoerc}, the factor $2\nu\|\cdots\|^2$, vanishes. Equal-order Galerkin
is therefore unusable at high $\Reyn$, and the stabilization exists precisely to
restore that lost control.

\section{Asymptotics I: $\tau_1$ saturates as $\Reyn\to\infty$}\label{sec:tausat}

The stabilization restores control through the parameter \eqref{eq:tau}. Its
behaviour as $\nu\to0$ is the crux.

\begin{proposition}[Saturation of $\tau_1$, consistent with paper \S\,\texttt{sec:Robustness}]\label{prop:tau}
Fix $h$, $|\bw|$, $\alpha_K$, $\sigma$ and let $\nu=\Reyn^{-1}\to0$. From the
definition of $\tauNS$ in \eqref{eq:tau}, the diffusive contribution
$c_1\nu/h^2$ drops relative to the convective $c_2|\bw|/h$ at a rate set by the
cell (mesh) Reynolds number $\Reyn_h:=|\bw|h/\nu$,
\begin{equation}
  \frac{\text{diffusive part of }\tauNS^{-1}}{\text{convective part}}
  =\frac{c_1\nu/h^2}{c_2|\bw|/h}
  =\frac{c_1}{c_2}\,\Reyn_h^{-1}\;\xrightarrow[\;\Reyn_h\to\infty\;]{}\;0,
  \label{eq:ratio}
\end{equation}
so that $\tauNS\to h/(c_2|\bw|)$ and hence
\begin{equation}
  \tau_1\;\longrightarrow\;\frac{1}{\alpha_K\,c_2|\bw|/h+\sigma}
  \;=\;
  \begin{cases}
    \dfrac{h}{\alpha_K c_2|\bw|}=\bigOh(h), & \text{convection-dominated }(\alpha_K c_2|\bw|/h\gg\sigma),\\[1.4ex]
    \dfrac{1}{\sigma}, & \text{reaction-dominated }(\sigma\gg\alpha_K c_2|\bw|/h),
  \end{cases}
  \label{eq:tausat}
\end{equation}
a finite, $\nu$-independent limit in both cases. These coincide with the paper's
robustness estimates $\tau_1\sim h/(\alpha\lVert\ba\rVert_{\infty,K})$ (dominant
convection) and $\tau_1\sim1/\sigma$ (dominant reaction). By contrast, in the
viscous limit $\Reyn_h\to0$ the paper gives $\tau_1\sim h^2/(\alpha_K\nu)\to0$.
\end{proposition}

For the momentum equation $\Reyn_h=|\bw|h/\nu$ is genuinely a Reynolds number, not
a P\'eclet number: what opposes advection here is the kinematic viscosity $\nu$
itself, the very $\nu$ of the global $\Reyn=UL/\nu$ (paper,
\texttt{eq:DimensionlessParameters}). It coincides with the cell P\'eclet number
of the generic advection--diffusion template (diffusivity an independent $k$) only
because the advected quantity is momentum.

\section{Asymptotics II: the stabilization residual does not vanish}\label{sec:rS}

It is convenient to read the discrete equation \eqref{eq:BS} as a single residual
functional. With $\Rop_{\bw}U_h=F-\Lop_{\bw}U_h$, the ASGS equation
$\BS=\LS|_{\boldsymbol\pi=0}$ is
\begin{equation}
  \rM(V):=\underbrace{\big[B(\bw,U_h,V)-L(V)\big]}_{\textstyle \rG(V)}
  \;+\;\underbrace{\sum_K\innerK{\Ladj_{\bw}V}{\tauK(\bw)\,\Rop_{\bw}U_h}}_{\textstyle \rS(V)}
  \;=\;0 ,
  \label{eq:resdecomp}
\end{equation}
(the OSGS variant replaces $\Rop U_h$ by $\Rop U_h-\boldsymbol\pi_h$). Hence at the
discrete root
\begin{equation}
  \rG(V)=-\,\rS(V)\qquad\forall V .
  \label{eq:cancel}
\end{equation}
The two functionals need not be individually small; only their sum must vanish.
Their size is governed by $\tau_1$. Since $\Rop_{\bw}U=F-\Lop_{\bw}U$ vanishes at
the exact solution, $\Rop_{\bw}U_h$ is, to leading order, the operator acting on
the error $\boldsymbol e_u=\buh-\bu$, $e_p=p_h-p$; as $\nu\to0$ it is carried by
the $\nu$-independent convective part $\alpha\,\ba\!\cdot\!\grad\boldsymbol e_u
+\alpha\grad e_p=\alpha X(U_h)|_{\text{err}}$, with $X(U_h):=\ba\!\cdot\!\grad\buh
+\grad p_h$ as in the paper's stability estimate. Combined with the saturated
$\tau_1=\bigOh(h/(\alpha_K|\bw|))$ of Proposition~\ref{prop:tau},

\begin{proposition}[Saturation of the stabilization residual]\label{prop:rS}
Let $S:=\|\rS\|$ be the dual norm of the stabilization functional at a near
solution. As $\nu\to0$ at fixed $h$, $S\to S_\infty=\bigOh(1)>0$
(convective limit), whereas $S\to0$ as $\nu\to\infty$ (viscous limit, where
$\tau_1\sim h^2/(\alpha_K\nu)\to0$). By \eqref{eq:cancel}, $\|\rG\|\to S_\infty$
too: \emph{both} pieces saturate to the same finite size.
\end{proposition}

This is the precise sense in which the stabilization \emph{restores} the control
that \eqref{eq:galcoerc} lost: the stabilized coercivity (paper,
\texttt{eq:StabilityEstimate}) adds the term
$\|\tau_1^{1/2}\alpha X(U_h)\|_h^2$ to the working norm, so the saturated $\tau_1$
now controls $\ba\!\cdot\!\grad\buh+\grad p_h$ that the vanishing viscous term no
longer could. The same $\tau_1$ that buys this control makes $\rS$ an
$\bigOh(1)$ contribution to the residual \eqref{eq:resdecomp}.

\section{The residual floor: cancellation digits demanded by the gate}\label{sec:floor}

The nonlinear solver drives the assembled vector $\rho_i=\rM(\boldsymbol\varphi_i)$
to zero and the gate declares success when
$\varepsilon_{\mathrm M}=\|\rho\|/D_{\mathrm M}\le\texttt{tol}=10^{-9}$, with
$D_{\mathrm M}=\bigOh(1)$ a $\tau$-free force-scale normalization. By
\eqref{eq:cancel} both $\rG$ and $\rS$ have the saturated size $S$, and the
equation residual is their remainder. Passing the gate requires the two large
pieces to agree to
\begin{equation}
  N_{\mathrm{digits}}\approx\log_{10}\frac{S}{\texttt{tol}\,D_{\mathrm M}}
  \label{eq:digits}
\end{equation}
decimal digits. As $\Reyn\to\infty$, $S\to S_\infty=\bigOh(1)$
(Prop.~\ref{prop:rS}) while $\texttt{tol},D_{\mathrm M}$ are fixed, so
$N_{\mathrm{digits}}$ rises to a large constant; the globalized Newton iteration
attains only $3$--$4$ digits of cancellation before the merit
$\Phi=\tfrac12\|\rM\|^2$ plateaus and the line search collapses, and
$\varepsilon_{\mathrm M}$ floors above \texttt{tol}. At low $\Reyn$, $S\to0$ and
$N_{\mathrm{digits}}\to0$: no cancellation is demanded and the residual reaches
machine precision. The floor is \emph{not} floating-point cancellation
($0.079-0.079$ is exact to $\sim10^{-17}$ in double precision) and \emph{not}
conditioning (it tracks $S$, not $\mathrm{cond}(J)$).

\paragraph{Numerical illustration.}
A representative convection-dominated cell ($\Reyn=10^5$, $N=10$, equal-order
$P_1/P_1$, ASGS) measured
$\|\rG\|\approx\|\rS\|\approx0.079$,
$\|\rM\|=\|\rG+\rS\|\approx2.4\times10^{-4}$,
$D_{\mathrm M}\approx2.85$, i.e. $\varepsilon_{\mathrm M}\approx8\times10^{-5}$
against $10^{-9}$. From \eqref{eq:digits},
$N_{\mathrm{digits}}=\log_{10}(0.079/(10^{-9}\!\cdot\!2.85))\approx7.4$: the gate
asks for $\sim7$--$8$ digit cancellation; the solver delivers $\sim3.5$. The
velocity error at this plateau is already optimal, $\|\boldsymbol e_u\|_0
=\bigOh(h^{k+1})$ -- as the paper's convergence theorem (paper,
\texttt{eq:ConvergenceResult}) guarantees -- so the leftover residual lies
\emph{below} the discretization error: the floor is a stopping-criterion
artifact, not an under-resolved solution.

\section{Contrast with pure Galerkin}\label{sec:galerkin}

Setting $\rS\equiv0$ in \eqref{eq:resdecomp} leaves $\rM=\rG$: no second large
term, so Newton drives $\rM$ to machine precision at \emph{any} $\Reyn$, with no
floor. But \eqref{eq:galcoerc} shows that residual belongs to a method with no
high-$\Reyn$ stability and (for equal order) no inf--sup control:

\begin{center}
\begin{tabular}{lcc}
\toprule
 & \textbf{Pure Galerkin} & \textbf{Stabilized (ASGS/OSGS)}\\
\midrule
Stabilization residual $\rS$        & absent                & present, $S\to S_\infty=\bigOh(1)$ \\
Velocity-gradient control as $\nu\to0$ & lost ($2\nu\|\cdots\|^2\to0$) & restored by $\|\tau_1^{1/2}\alpha X(U_h)\|_h^2$\\
Equal-order inf--sup (LBB)          & violated              & restored (Lemma, paper \texttt{lemma:Stability})\\
High cell-$\Reyn_h$ oscillations    & present               & damped (SUPG-type upwinding)\\
High-$\Reyn$ absolute residual floor& none ($\to$ machine $\epsilon$) & $\varepsilon_{\mathrm M}$ floors $\gg\texttt{tol}$\\
Usable at high $\Reyn$, equal order & \textbf{no}           & \textbf{yes}\\
\bottomrule
\end{tabular}
\end{center}

The term that makes equal-order interpolation stable -- $\sum_K\langle\Ladj V,
\tauK\Lop U\rangle_K$, equivalently the control $\|\tau_1^{1/2}\alpha X(U_h)\|_h^2$
in the paper's coercivity bound -- is exactly the term whose $\tau_1$-saturation
floors the residual. The floor is the \emph{price} of consistent stabilization,
visible only under a tight absolute gate. Galerkin avoids the price by not buying
stability; it is a foil that isolates the mechanism, not a usable alternative.

\section{Scope: the two layers, both consistent with the paper}\label{sec:scope}

\paragraph{Layer 1 (this note) is generic, and is a high-$\Reyn$ phenomenon.}
Nothing in \S\ref{sec:galcoerc}--\ref{sec:floor} uses $\sigma$ essentially:
putting $\bsig\equiv0$ (i.e. $\alpha\equiv1$, the standard Navier--Stokes system
recovered by the paper at $\alpha\equiv1$) leaves the argument intact. The
high-$\Reyn$ residual floor is therefore generic; it is invisible in ordinary
practice only because ordinary practice stops on a \emph{relative} residual and
never drives a scale-free absolute residual to $10^{-9}$ on the finest
manufactured-solution mesh. The paper's Galerkin-coercivity remark
(\S\ref{sec:galcoerc}) is the continuous-level statement of the same loss of
control. Empirically (\S\ref{sec:experiment}) this floor sets in only at very high
$\Reyn$: at $\Reyn=10^5$ the near-Navier--Stokes cell $\alpha_0=0.9$ still reaches
the $10^{-9}$ gate with optimal $\bigOh(h^{k+1})$ velocity error; the regular
harness first shows it at $\Reyn=10^6$.

\paragraph{Layer 2 (low porosity) is the paper's ubiquitous $1/\alpha_0$ degradation,
driven by reaction \emph{magnitude}.}
The reason the \emph{low-porosity} cells fold -- already at $\Reyn=10^5$, where
Layer 1 is dormant -- is a \emph{separate} effect that the paper's robustness
analysis predicts. Its convergence estimates carry an overall factor $1/\alpha_0$
in every regime (paper, \texttt{eq:ConvergenceResultDominantViscosity} and the
following remark: ``this linear drop in accuracy with decreasing minimal porosity
is ubiquitous over the space of physical parameters\dots\ because the terms
involving derivatives of the velocity are multiplied by the porosity''); the
working norm itself weights the velocity gradient by $\alpha$ (paper,
\texttt{eq:SigmaAlpha}, $\sigtil$, and the triple norm); and in the
reaction-dominated regime $\Damk_h\to\infty$ (paper, \S\,\texttt{sec:DominantReaction})
the parameter degenerates to $\tau_1\sim1/\sigma$ with the weakened control
$\sigtil$. At low $\alpha_0$ the linear Darcy drag $a(\alpha)\propto((1-\alpha)/\alpha)^2$
makes $\sigma$ -- and hence $\Damk$ -- large, pushing the cell into exactly this
regime. \emph{The driver is the reaction magnitude, not its nonlinearity}: the
controlled experiment of \S\ref{sec:experiment} shows that removing the
$b(\alpha)|\bu|$ term (so $\partial\sigma/\partial\bu=0$) at matched $\Damk$ folds
identically, so the $|\bu|$-coupling basin-shrink is \emph{not} the cause.

\paragraph{Numerical confirmation.}\label{sec:experiment}
Holding $\Reyn=10^5$ fixed (ASGS, $P_1/P_1$, $N=10,20,40$) and sweeping the minimum
porosity gives a clean onset: $\alpha_0=0.9$ and $0.5$ \textbf{converge} to the
$10^{-9}$ gate with optimal rate $\approx2.2$, while $\alpha_0=0.2$ and $0.1$
\textbf{fold} (large, pre-asymptotic $\|\boldsymbol e_u\|$). A matched-$\Damk$
\emph{linear}-reaction control at $\alpha_0=0.1$ ($b\equiv0$, $a$ rescaled to keep
$\Damk_h$ fixed) also folds, with essentially the same errors as the full
Forchheimer case. The two controls jointly locate the difficulty: it is the
low-porosity reaction \emph{magnitude} (Layer 2), not the high-$\Reyn$
$\tau_1$-saturation floor (Layer 1, dormant at $\Reyn=10^5$) and not the
Forchheimer $|\bu|$ nonlinearity.

\begin{center}
\fbox{\parbox{0.93\linewidth}{\centering
\textbf{Layer 1} ($\tau_1$ saturation, this note): a high-$\Reyn$ residual floor
under a tight absolute gate -- \emph{generic}, Navier--Stokes included, the
solver shadow of the paper's ``almost no control over the velocity gradient'';
empirically dormant until $\Reyn\!\gtrsim\!10^6$.\\[3pt]
\textbf{Layer 2} (low $\alpha_0$): the paper's ubiquitous $1/\alpha_0$ accuracy
degradation and reaction-dominated regime $\tau_1\sim1/\sigma$, driven by reaction
\emph{magnitude} (confirmed: the $|\bu|$ nonlinearity is \emph{not} the cause) --
\emph{porous-specific}, why low-porosity cells fold already at $\Reyn=10^5$.}}
\end{center}

The two are consistent: the paper proves the \emph{solution} stays optimal in the
triple norm (paper, \texttt{eq:ConvergenceResult}); this note explains why the
\emph{algebraic} residual nonetheless cannot reach a tight absolute gate at high
$\Reyn$, and why low porosity makes the nonlinear solve harder still.

\end{document}
